% GENERATED FILE. Edit canonical lessons, then run npm run generate:books.
% canonical-source-hash: fnv1a64:9f5d67531789831e
% canonical-lessons: ZH-C335-yuanxiao, ZH-C335-denglong, ZH-C335-cai, ZH-C335-miyu, ZH-C335-bianpao, ZH-R335-review-330-335

\chapter{Lanterns and Riddles}
\label{ch:zh-lanterns-and-riddles}
\hlchaptermodality{335}

\begin{chapteropening}
\textbf{By the end of this chapter:} \emph{I can talk about the Lantern Festival, lanterns, riddles and firecrackers, and say I guessed.}

Complete the last lesson of chapter 335: I can talk about the Lantern Festival, lanterns, riddles and firecrackers, and say I guessed.
\end{chapteropening}

\section[Yuánxiāo]{Yuánxiāo — the Lantern Festival}

\label{lesson:ZH-C335-yuanxiao}



\begin{quote}
Before the new one: say the Mandarin for the Spring Festival couplets, then the Mandarin for flourishing in business.
\end{quote}

\subsection*{You'll want to know: Yuánxiāo}
\emph{Yuánxiāo} — \textquotedblleft{}the Lantern Festival\textquotedblright{}.

\begin{quote}
\emph{Yuánxiāo de wǎnshang wǒmen chūqù wán} — On the night of the lantern festival we go out and have fun.
\end{quote}

Its characters come later; for now it is a word for the ear.

\begin{cousinweb}
\emph{Yuánxiāo} is built from \zh{元} 'first' and \emph{xiāo} 'night', both inherited from Middle Chinese, \emph{ngjwon} and \emph{sjew} (Guangyun, in Baxter's transcription), which glosses the first 'great' and 'the beginning' and the second 'night': the first full-moon night of the year.
\end{cousinweb}

\begin{grammarlens}[title={What you've built}]
One more word for telling what happened.
\end{grammarlens}

\subsection*{Your turn}
\begin{itemize}
\raggedright
  \item \emph{Say it:} \emph{Yuánxiāo}
  \item \emph{Say it:} \emph{Yuánxiāo}, once more
  \item \emph{Say it:} \emph{Yuánxiāo de wǎnshang wǒmen chūqù wán} — On the night of the lantern festival we go out and have fun.
  \item \emph{Recall:} say the Mandarin for the Spring Festival couplets, then the Mandarin for flourishing in business, then say \emph{Yuánxiāo} again
\end{itemize}

\begin{tcolorbox}[breakable,colback=teal!4,colframe=teal!35!black,title={Before you move on}]
What does Yuánxiāo mean? (\textquotedblleft{}The Lantern Festival\textquotedblright{}.) Say it once more.
\end{tcolorbox}
\section[dēnglóng]{dēnglóng — lantern}

\label{lesson:ZH-C335-denglong}



\begin{quote}
Before the new one: say the Mandarin for flourishing in business, then the Mandarin for the Lantern Festival.
\end{quote}

\subsection*{You'll want to know: dēnglóng}
\emph{dēnglóng} — \textquotedblleft{}lantern\textquotedblright{}.

\begin{quote}
\emph{ménkǒu yǒu liǎng ge hóng dēnglóng} — There are two red lanterns at the door.
\end{quote}

Its characters come later; for now it is a word for the ear.

\begin{cousinweb}
\emph{dēnglóng} is built from \emph{dēng} 'a lamp' and \emph{lóng} 'a basket', both inherited from Middle Chinese, \emph{tong} and \emph{luwng} (Guangyun, in Baxter's transcription), which glosses the first 'a lamp' and quotes an older book for the second, a basket for a fire.
\end{cousinweb}

\begin{grammarlens}[title={What you've built}]
One more word for telling what happened.
\end{grammarlens}

\subsection*{Your turn}
\begin{itemize}
\raggedright
  \item \emph{Say it:} \emph{dēnglóng}
  \item \emph{Say it:} \emph{dēnglóng}, once more
  \item \emph{Say it:} \emph{ménkǒu yǒu liǎng ge hóng dēnglóng} — There are two red lanterns at the door.
  \item \emph{Recall:} say the Mandarin for flourishing in business, then the Mandarin for the Lantern Festival, then say \emph{dēnglóng} again
\end{itemize}

\begin{tcolorbox}[breakable,colback=teal!4,colframe=teal!35!black,title={Before you move on}]
What does dēnglóng mean? (\textquotedblleft{}Lantern\textquotedblright{}.) Say it once more.
\end{tcolorbox}
\section[cāi]{cāi — to guess}

\label{lesson:ZH-C335-cai}



\begin{quote}
Before the new one: say the Mandarin for the Lantern Festival, then the Mandarin for a paper lantern.
\end{quote}

\subsection*{You'll want to know: cāi}
\emph{cāi} — \textquotedblleft{}to guess\textquotedblright{}.

\begin{quote}
\emph{wǒ cāi le sān cì dōu bú duì} — I guessed three times and got it wrong each time.
\end{quote}

Its characters come later; for now it is a word for the ear.

\begin{cousinweb}
\emph{cāi} is inherited from Middle Chinese \emph{tshoj} (Guangyun, in Baxter's transcription), which glosses it 'to suspect' and 'to resent'. The sense 'to guess' is not given there.
\end{cousinweb}

\begin{grammarlens}[title={What you've built}]
One more word for telling what happened.
\end{grammarlens}

\subsection*{Your turn}
\begin{itemize}
\raggedright
  \item \emph{Say it:} \emph{cāi}
  \item \emph{Say it:} \emph{cāi}, once more
  \item \emph{Say it:} \emph{wǒ cāi le sān cì dōu bú duì} — I guessed three times and got it wrong each time.
  \item \emph{Recall:} say the Mandarin for the Lantern Festival, then the Mandarin for a paper lantern, then say \emph{cāi} again
\end{itemize}

\begin{tcolorbox}[breakable,colback=teal!4,colframe=teal!35!black,title={Before you move on}]
What does cāi mean? (\textquotedblleft{}To guess\textquotedblright{}.) Say it once more.
\end{tcolorbox}
\section[míyǔ]{míyǔ — riddle}

\label{lesson:ZH-C335-miyu}



\begin{quote}
Before the new one: say the Mandarin for a paper lantern, then the Mandarin for to guess an answer.
\end{quote}

\subsection*{You'll want to know: míyǔ}
\emph{míyǔ} — \textquotedblleft{}riddle\textquotedblright{}.

\begin{quote}
\emph{Yuánxiāo de shíhou wǒmen cāi míyǔ} — At the lantern festival we guess riddles.
\end{quote}

Its characters come later; for now it is a word for the ear.

\begin{cousinweb}
\emph{míyǔ} is built from \emph{mí} and \zh{语} 'speech', both inherited from Middle Chinese, \emph{mejH} and \emph{ngjoX} (Guangyun, in Baxter's transcription), which glosses the first 'veiled words'.
\end{cousinweb}

\begin{grammarlens}[title={What you've built}]
One more word for telling what happened.
\end{grammarlens}

\subsection*{Your turn}
\begin{itemize}
\raggedright
  \item \emph{Say it:} \emph{míyǔ}
  \item \emph{Say it:} \emph{míyǔ}, once more
  \item \emph{Say it:} \emph{Yuánxiāo de shíhou wǒmen cāi míyǔ} — At the lantern festival we guess riddles.
  \item \emph{Recall:} say the Mandarin for a paper lantern, then the Mandarin for to guess an answer, then say \emph{míyǔ} again
\end{itemize}

\begin{tcolorbox}[breakable,colback=teal!4,colframe=teal!35!black,title={Before you move on}]
What does míyǔ mean? (\textquotedblleft{}Riddle\textquotedblright{}.) Say it once more.
\end{tcolorbox}
\section[biānpào]{biānpào — firecrackers}

\label{lesson:ZH-C335-bianpao}



\begin{quote}
Before the new one: say the Mandarin for to guess an answer, then the Mandarin for a riddle.
\end{quote}

\subsection*{You'll want to know: biānpào}
\emph{biānpào} — \textquotedblleft{}firecrackers\textquotedblright{}.

\begin{quote}
\emph{Chúxī wǎnshang shí'èr diǎn wàibian de rén dōu zài fàng biānpào} — At midnight on the last night of the year everyone outside is letting off firecrackers.
\end{quote}

Its characters come later; for now it is a word for the ear.

\begin{cousinweb}
\emph{biānpào} is built from \emph{biān} 'a whip' and \emph{pào}, both inherited from Middle Chinese, \emph{pjien} and \emph{phaewH} (Guangyun, in Baxter's transcription), which glosses the first 'a horse whip' and the second 'the look of something scorched'.
\end{cousinweb}

\begin{grammarlens}[title={What you've built}]
One more word for telling what happened.
\end{grammarlens}

\subsection*{Your turn}
\begin{itemize}
\raggedright
  \item \emph{Say it:} \emph{biānpào}
  \item \emph{Say it:} \emph{biānpào}, once more
  \item \emph{Say it:} \emph{Chúxī wǎnshang shí'èr diǎn wàibian de rén dōu zài fàng biānpào} — At midnight on the last night of the year everyone outside is letting off firecrackers.
  \item \emph{Recall:} say the Mandarin for to guess an answer, then the Mandarin for a riddle, then say \emph{biānpào} again
\end{itemize}

\begin{tcolorbox}[breakable,colback=teal!4,colframe=teal!35!black,title={Before you move on}]
What does biānpào mean? (\textquotedblleft{}Firecrackers\textquotedblright{}.) Say it once more.
\end{tcolorbox}
\section[(dialogue)]{Looking back: from The King's Son to Lanterns and Riddles}

\label{lesson:ZH-R335-review-330-335}



\begin{quote}
Say the Mandarin for to inherit the family business and a prince.
\end{quote}

\subsection*{Your turn}
\begin{itemize}
\raggedright
  \item \emph{Say it:} to inherit the family business, a prince, a king in a story, a king's wife, a kingdom ruled by a king
  \item \emph{Say it:} a monster in a story, a Daoist immortal, magic in a fairy tale, a precious treasure, to hide away
  \item \emph{Say it:} the nerve to face danger, a ghost in a story, to be haunted by a ghost, how absurd, the human world below
  \item \emph{Say it:} the Spring Festival Gala on TV, lunar New Year's Eve, the New Year's Eve dinner, to keep the New Year vigil, to pay a New Year call
  \item \emph{Say it:} the first lunar month, a red envelope of gift money, to paste up on a door, the Spring Festival couplets, flourishing in business
  \item \emph{Say it:} the Lantern Festival, a paper lantern, to guess an answer, a riddle, a string of firecrackers
  \item \emph{Read it:} \zh{王} \zh{鬼} \zh{红} — each one aloud, then a word that holds it
\end{itemize}

\begin{tcolorbox}[breakable,colback=teal!4,colframe=teal!35!black,title={Before you move on}]
To inherit the family business? (\emph{jìchéng}.) A prince? (\textbf{\zh{王子}}, \emph{wángzǐ}.) A king in a story? (\textbf{\zh{国王}}, \emph{guówáng}.) A king's wife? (\textbf{\zh{王后}}, \emph{wánghòu}.) A kingdom ruled by a king? (\textbf{\zh{王国}}, \emph{wángguó}.) A monster in a story? (\emph{yāoguài}.) A Daoist immortal? (\emph{shénxiān}.) Magic in a fairy tale? (\emph{mófǎ}.) A precious treasure? (\emph{bǎowù}.) To hide away? (\emph{cáng}.) The nerve to face danger? (\emph{dǎnzi}.) A ghost in a story? (\textbf{\zh{鬼}}, \emph{guǐ}.) To be haunted by a ghost? (\emph{nàoguǐ}.) How absurd? (\textbf{\zh{见鬼}}, \emph{jiànguǐ}.) The human world below? (\textbf{\zh{人间}}, \emph{rénjiān}.) The Spring Festival Gala on TV? (\textbf{\zh{春晚}}, \emph{Chūnwǎn}.) Lunar New Year's Eve? (\emph{Chúxī}.) The New Year's Eve dinner? (\emph{niányèfàn}.) To keep the New Year vigil? (\emph{shǒusuì}.) To pay a New Year call? (\emph{bàinián}.) The first lunar month? (\textbf{\zh{正月}}, \emph{zhēngyuè}.) A red envelope of gift money? (\textbf{\zh{红包}}, \emph{hóngbāo}.) To paste up on a door? (\emph{tiē}.) The Spring Festival couplets? (\emph{chūnlián}.) Flourishing in business? (\textbf{\zh{红火}}, \emph{hónghuǒ}.) The Lantern Festival? (\emph{Yuánxiāo}.) A paper lantern? (\emph{dēnglóng}.) To guess an answer? (\emph{cāi}.) A riddle? (\emph{míyǔ}.) A string of firecrackers? (\emph{biānpào}.)
\end{tcolorbox}
